\begin{abstract}
Conventional Hyperspectral Imaging (HSI) is based on acquiring either reflectance, emission, or excitation profiles from each pixel of an image.
Emission-based HSI can be expanded into a fourth dimension by varying the excitation wavelength, which enables an improved discrimination of materials that contain mixed fluorophores.
Multi-excitation HSI (ME-HSI) creates a 4D dataset encompassing signal intensity across two spatial and two spectral dimensions, the latter being the emission and excitation wavelengths.
The processing of ME-HSI datasets presents notable challenges due to the numerous spectral layers per pixel, introducing considerable redundancy and computational burden.
Existing band selection methods, developed for 3D HSI, fail to capture the nonlinear cross-excitation correlations present in ME-HSI.
To fill this gap, we developed a deep learning framework comprising three stages: (1) a 3D convolutional autoencoder with parallel excitation branches that learns a unified latent representation of ME-HSI data, (2) a perturbation-based attribution method that traces reconstruction sensitivity to individual excitation-emission wavelength pairs, and (3) Maximum Marginal Relevance (MMR) selection that balances informativeness with spectral diversity.
The framework is evaluated on three datasets covering distinct ME-HSI regimes: a four-class lichen autofluorescence dataset (3{,}072 configurations sweep), a three-class collagen sponges--collagen dataset that tests transfer to a chemically distinct fluorophore family, and a fully unlabeled droplet dataset that tests the unsupervised claim end-to-end.
On lichens, 95.2\% accuracy is achieved with 80 of 192 bands (58\% data reduction, 7.0 points above the 88.2\% baseline), and 89.4\% is maintained at only 9 bands (95\% reduction).
On collagen sponges--collagen, the same pipeline---unchanged---reaches 85.6\% accuracy at 30 of 158 bands (81\% reduction, 5.8 points above baseline); a ten-classifier evaluation confirms the selected bands are intrinsically informative rather than tuned to a single classifier.
On the unlabeled droplet dataset, the framework recovers the three spectral types identified by full-spectrum hierarchical clustering at 96.4\% per-pixel cross-validated accuracy using only 10 of 214 bands, outperforming or matching eight band-selection baselines including BS-Net-FC and a supervised Sparse-LASSO control.
The framework facilitates practical applications of ME-HSI by reducing acquisition complexity while improving discriminative power, with results that hold across supervised and fully blind settings.
\end{abstract}